% POLYDIM_DEST
% destino: polydim/docs/
% nombre:  polydim_paper_arxiv.tex
% autor:   ai.mpat.agt@gmail.com
% fecha:   2026-06-26
% tarea:   TASK_LATEX
% base:    POLYDIM_PAPER_COMPLETO_V3.md
% target:  arXiv cs.PL + cs.AI
% compilar: pdflatex polydim_paper_arxiv.tex && bibtex polydim_paper_arxiv && pdflatex x2

\documentclass[10pt,twocolumn]{article}

%% ============================================================
%% Packages
%% ============================================================
\usepackage[utf8]{inputenc}
\usepackage[T1]{fontenc}
\usepackage[margin=1in]{geometry}
\usepackage{amsmath,amssymb,amsthm}
\usepackage{mathtools}
\usepackage{bm}           % bold math
\usepackage{microtype}    % better typography
\usepackage{hyperref}
\usepackage{cleveref}
\usepackage{booktabs}     % professional tables
\usepackage{listings}     % code listings
\usepackage{xcolor}
\usepackage{graphicx}
\usepackage{enumitem}
\usepackage{authblk}
\usepackage{natbib}       % bibliography

%% ============================================================
%% Custom colors
%% ============================================================
\definecolor{polyblue}{RGB}{30,80,160}
\definecolor{codegray}{RGB}{248,248,248}
\definecolor{codegreen}{RGB}{0,128,0}
\definecolor{codered}{RGB}{160,0,0}
\definecolor{codepurple}{RGB}{120,0,120}

%% ============================================================
%% Listing styles
%% ============================================================
\lstdefinelanguage{polydim}{
  keywords={COMPOSE,MIX,FIXPOINT,COMPILE,RENDER,EXPORT,ATTEND,RECUR,BIND,
            position,transform,output,bind},
  keywordstyle=\color{polyblue}\bfseries,
  ndkeywords={DIM_FLUTTER,DIM_RUST,DIM_WASM,DIM_SQL,DIM_PYTHON,
              DIM_GRAPH,DIM_VECTOR,DIM_TIME,DIM_ERROR,DIM_META},
  ndkeywordstyle=\color{codepurple},
  sensitive=true,
  comment=[l]{--},
  commentstyle=\color{codegreen}\itshape,
  stringstyle=\color{codered},
  morestring=[b]",
}

\lstset{
  basicstyle=\ttfamily\small,
  backgroundcolor=\color{codegray},
  frame=single,
  framesep=3pt,
  numbers=none,
  breaklines=true,
  captionpos=b,
  tabsize=2,
  showstringspaces=false,
}

%% ============================================================
%% Theorem environments
%% ============================================================
\theoremstyle{plain}
\newtheorem{theorem}{Theorem}[section]
\newtheorem{corollary}[theorem]{Corollary}
\newtheorem{lemma}[theorem]{Lemma}
\newtheorem{proposition}[theorem]{Proposition}

\theoremstyle{definition}
\newtheorem{definition}[theorem]{Definition}
\newtheorem{example}[theorem]{Example}

\theoremstyle{remark}
\newtheorem{remark}[theorem]{Remark}

%% ============================================================
%% Math macros
%% ============================================================
\newcommand{\RR}{\mathbb{R}}
\newcommand{\NN}{\mathbb{N}}
\newcommand{\GeoID}{\mathsf{GEO\_ID}}
\newcommand{\Compose}{\mathsf{COMPOSE}}
\newcommand{\Mix}{\mathsf{MIX}}
\newcommand{\Fixpoint}{\mathsf{FIXPOINT}}
\newcommand{\Project}{\mathsf{PROJECT}}
\newcommand{\Attend}{\mathsf{ATTEND}}
\newcommand{\Recur}{\mathsf{RECUR}}
\newcommand{\Align}{\mathsf{ALIGN}}
\newcommand{\DimX}{\mathsf{DIM}_X}
\newcommand{\ProxX}{\Pi_X}
\newcommand{\EmbX}{\mathrm{embed}_X}
\newcommand{\ProjX}{\mathrm{proj}_X}
\newcommand{\Ob}{\mathrm{Ob}}
\newcommand{\Hom}{\mathrm{Hom}}
\newcommand{\norm}[1]{\left\|#1\right\|}
\newcommand{\lfp}{\mathrm{lfp}}
\newcommand{\dom}{\mathrm{Dom}}

%% ============================================================
%% Title, authors
%% ============================================================
\title{\textbf{POLYDIM: A Transformer-Native Algebraic Programming Language\\
for High-Dimensional Latent Computing}}

\author[1]{POLYDIM Research Group}
\affil[1]{\texttt{ai.mpat.agt@gmail.com}}
\date{2026}

%% ============================================================
\begin{document}
%% ============================================================

\maketitle

%% ============================================================
\begin{abstract}
Modern artificial intelligence architectures---Transformers and State Space
Models (SSMs)---perform computation through high-dimensional geometric
transformations in latent spaces of thousands to tens of thousands of
dimensions. Yet these architectures are programmed and forced to communicate
via discrete sequential symbol streams, collapsing the rich geometry of their
latent states to a one-dimensional text channel. We define this as the
\emph{Impedance Mismatch}: a paradigm gap requiring constant
serialization-deserialization cycles, producing semantic loss with every
AI-to-AI interaction.

We propose \textsc{Polydim}, a programming language whose fundamental unit of
computation is a geometric transformation $T\colon \RR^N \to \RR^N$. A closed
algebraic core of four operators---$\Compose$, $\Mix$, $\Fixpoint$, and
$\Project$---enables native latent-space execution, zero-serialization AI-to-AI
communication via cross-attention semantics, and a dimensional type system where
types emerge from orthogonal projections rather than declarations. We establish
four algebraic theorems: associativity of composition, linearity of
superposition, \emph{functoriality of projection} (Theorem~3, the central
contribution), and uniqueness of fixed points. The language specification
incorporates Low-Rank Adaptation (LoRA) serialization for practical
distribution, and Homotopy Type Theory~(HoTT) for model-agnostic object
identity. A Python bootstrap achieves $29/29$ tests and a $64.5\%$ semantic
improvement over deterministic baselines.
\end{abstract}

%% ============================================================
\section{Introduction}
\label{sec:intro}
%% ============================================================

\subsection{The Impedance Mismatch Problem}

Consider a large language model with an internal representation space of
dimension $d = 12{,}288$ (GPT-4 scale). When two such models communicate, the
standard protocol requires:
(1)~collapsing the high-dimensional state to approximately 100 sequential tokens
(serialize); (2)~transmitting the tokens; (3)~re-expanding from tokens to the
high-dimensional state (deserialize); and (4)~re-computing the geometric context
from the sequential text. Two irreversible projections, two semantic losses, no
guarantees of fidelity at the receiving end.

\paragraph{The Gesture vs.~Emoji Analogy.} An AI operating natively works with
a ``gesture''---a continuous, high-fidelity manifold of meaning in $1{,}000$+
dimensions. Forcing text communication collapses that gesture to a flat
two-dimensional ``emoji'': functionally recognizable, inherently lossy.
\textsc{Polydim} allows AIs to exchange the full gesture directly as tensor
transformations.

The historical parallel is precise:
\begin{quote}
\textit{1950s: CPUs already performed binary arithmetic.\\
$\rightarrow$ Assembly formalized how to program them deliberately.\\[4pt]
2020s: Transformers already operate in multidimensional geometry.\\
$\rightarrow$ \textsc{Polydim} formalizes how to program them deliberately.}
\end{quote}

\subsection{Contributions}

\begin{enumerate}[leftmargin=*,label=\arabic*.]
\item \textbf{Language specification:} formally specified with BNF grammar,
      operational semantics, and denotational semantics.
\item \textbf{Two-layer architecture:} invariant algebraic core separating
      architecture-agnostic operations from transformer/SSM-specific primitives.
\item \textbf{Dimensional type system:} types emerge from projection, formalized
      as a functor $F\colon \mathcal{G} \to \mathcal{E}$ (Theorem~\ref{thm:functor}).
\item \textbf{Zero-serialization AI$\leftrightarrow$AI protocol:} the $\Align$
      protocol via Procrustes eliminates text serialization.
\item \textbf{HoTT-based identity:} $\GeoID$ as a homotopy equivalence class,
      enabling model-agnostic object identity.
\item \textbf{Python bootstrap:} $29/29$ tests, align\_score $= 0.9993$,
      $+64.5\%$ semantic improvement over deterministic baselines.
\end{enumerate}

%% ============================================================
\section{Background}
\label{sec:background}
%% ============================================================

\subsection{Transformers as Geometric Computation}

The central operation of a Transformer~\citep{vaswani2017} is:
\[
  \mathrm{Attention}(Q, K, V)
  = \mathrm{softmax}\!\left(\frac{QK^\top}{\sqrt{d}}\right) V
\]
This is $T\colon \RR^d \to \RR^d$---a simultaneous transformation of the full
representation space. Every transformer layer is $s_{\mathrm{out}} = T_L \circ
\cdots \circ T_1(s_{\mathrm{in}})$, i.e., $\Compose$ applied $L$
times---\textsc{Polydim}'s fundamental operation, applied natively by every
production transformer.

\subsection{State Space Models}

Mamba~\citep{gu2023mamba} and related SSMs provide:
$h(t+1) = Ah(t) + Bx(t)$, $y(t) = Ch(t)$. The matrices $A, B, C$ constitute
the program; the state $h$ is the computation. A \texttt{.polydim} file can
literally be these matrices serialized as a program.

\subsection{VSA and Quasi-Orthogonality}

Vector Symbolic Architectures~\citep{plate1995,kanerva2009} demonstrate that
algebraic operations over high-dimensional random vectors preserve semantic
structure. A critical property: for large $N$, two random vectors in $\RR^N$
are quasi-orthogonal: $\mathbb{E}[\cos(u,v)] = 0$, $\mathrm{Var}[\cos(u,v)] =
1/N$. For $N = 10{,}000$: $\sigma \approx 0.01$. This makes two simultaneous
branches of $\Mix$ coexist without destructive interference
(Section~\ref{sec:language}).

\subsection{Category Theory Foundations}

A category $\mathcal{C}$ consists of objects, morphisms, identity morphisms, and
associative composition. A \emph{functor} $F\colon \mathcal{C} \to \mathcal{D}$
preserves identity ($F(\mathrm{id}_A) = \mathrm{id}_{F(A)}$) and composition
($F(g \circ f) = F(g) \circ F(f)$). These structures underlie
Theorem~\ref{thm:functor}.

%% ============================================================
\section{POLYDIM Language Design}
\label{sec:language}
%% ============================================================

\subsection{Core Principles}

\begin{description}[leftmargin=1.5em]
\item[P1] No instructions, only transformations. The fundamental unit is
  $T\colon \RR^N \to \RR^N$.
\item[P2] No variables, only positions. A position $P = (g, S)$ in $\RR^N$;
  its type emerges from observation.
\item[P3] No sequential loops, only convergence. $\Fixpoint(T, \varepsilon)$
  replaces iteration.
\item[P4] No binary branching, only superposition. $\Mix(\alpha, T_1, \beta,
  T_2)$ replaces \texttt{if/else}.
\end{description}

\subsection{Formal Definitions and Grammar}
\label{sec:defs}

\begin{definition}[State]
$S = (V, D, A)$ where $V \in \RR^N$, $D \subseteq \{\mathsf{DIM}_0, \ldots,
\mathsf{DIM}_8\}$, and $A\colon D \to [0,1]$.
\end{definition}

\begin{definition}[Position]
$P = (g, S)$ where $g \in \RR^N$ is the $\GeoID$ (invariant under all
admissible transformations) and $S$ is the mutable state.
\end{definition}

\begin{definition}[Admissible Transformation]
$T\colon \RR^N \to \RR^N$ is admissible if: (i)~$T$ is measurable;
(ii)~$\norm{T(g) - g} < \varepsilon_g$ for all $\GeoID$s $g$;
(iii)~$T = W_0 + UV^\top$ with $U, V \in \RR^{N \times r}$, $r \ll N$.
\end{definition}

\begin{definition}[POLYDIM Program]
$\Pi = T_n \circ T_{n-1} \circ \cdots \circ T_1$, a finite composition of
admissible transformations.
\end{definition}

\paragraph{Formal Grammar (BNF).}

\begin{lstlisting}[language=polydim,caption={Abstract BNF grammar of POLYDIM.}]
<program>   ::= <transform>
<transform> ::=
    COMPOSE ( <transform> , <transform> )
  | MIX ( <weight> , <transform> ,
              <weight> , <transform> )
  | FIXPOINT ( <transform> , <epsilon> )
  | COMPILE ( <transform> , <executor> )
  | RENDER  ( <transform> , <executor> )
  | EXPORT  ( <transform> , <executor> )
  | ATTEND ( <matrix> , <matrix> , <matrix> )
  | RECUR  ( <matrix> , <matrix> , <matrix> )
  | <position-ref>
<executor>  ::=
    DIM_FLUTTER | DIM_RUST | DIM_SQL
  | DIM_PYTHON  | DIM_WASM | DIM_GRAPH
  | <executor> ∩ <executor>
<weight>    ::= float ∈ [0.0, 1.0]
<epsilon>   ::= float ∈ (0.0, 1.0)
\end{lstlisting}

The grammar excludes named variables~(P2), sequential loops~(P3), binary
branching~(P4), and declared types---defining absences, not restrictions.

\subsection{Two-Layer Architecture}

\textbf{Algebraic Layer} (invariant):
\begin{align*}
\Compose(T_1, T_2)      &= T_2 \circ T_1 \\
\Mix(\alpha, T_1, \beta, T_2) &= \alpha T_1 + \beta T_2 \\
\Fixpoint(T, \varepsilon) &= T^n(v) \text{ until } \norm{s_{k+1} - s_k} < \varepsilon \\
\Project(T, E)           &= F_E(T(s))
\end{align*}

\textbf{Implementation Layer} (architecture-specific):
\begin{align*}
\Attend(W_Q, W_K, W_V, s) &= \mathrm{softmax}\!\left(\tfrac{(sW_Q)(sW_K)^\top}{\sqrt{d}}\right)(sW_V) \\
\Recur(A, B, C, h, x)    &= Ah + Bx,\quad y = Ch
\end{align*}

\subsection{Big-Step Operational Semantics}

We write $\langle s, T \rangle \Downarrow s'$ for ``state $s$ evaluates under
$T$ to $s'$'':

\begin{gather*}
\frac{\langle s, T_1 \rangle \Downarrow s_1 \quad \langle s_1, T_2 \rangle \Downarrow s_2}
     {\langle s, \Compose(T_1, T_2) \rangle \Downarrow s_2}
\tag{E-Compose}\\[6pt]
\frac{}
     {\langle s, \Mix(\alpha, T_1, \beta, T_2) \rangle \Downarrow \alpha T_1(s) + \beta T_2(s)}
\tag{E-Mix}\\[6pt]
\frac{\exists n\colon \norm{T^n(s) - T^{n-1}(s)} < \varepsilon \quad s^* = T^n(s)}
     {\langle s, \Fixpoint(T, \varepsilon) \rangle \Downarrow s^*}
\tag{E-Fix}\\[6pt]
\frac{F_E\colon \mathcal{G} \to \mathcal{E}}
     {\langle s, \Project(T, E) \rangle \Downarrow F_E(T(s))}
\tag{E-Proj}
\end{gather*}

\subsection{Formal Theorems}
\label{sec:theorems}

\begin{theorem}[Associativity of $\Compose$]
\label{thm:assoc}
For all admissible $T_1, T_2, T_3$:
\[
  \Compose(\Compose(T_3, T_2), T_1) = \Compose(T_3, \Compose(T_2, T_1)).
\]
\end{theorem}

\begin{proof}
Follows from standard associativity of function composition on $\RR^N$. \qed
\end{proof}

\begin{remark}
Non-commutativity ($\Compose(T_2, T_1) \neq \Compose(T_1, T_2)$ in general)
encodes causal order topologically in the algebra rather than positionally in an
instruction sequence.
\end{remark}

\begin{theorem}[Linearity of $\Mix$]
\label{thm:linear}
If $T_1, T_2 \in \mathrm{Lin}(\RR^N)$, then
$\Mix(\alpha, T_1, \beta, T_2) \in \mathrm{Lin}(\RR^N)$.
\end{theorem}

\begin{proof}
A linear combination of linear maps is linear. \qed
\end{proof}

\begin{remark}[VSA Foundation]
For $N = 10{,}000$, $\sigma[\cos(u,v)] \approx 0.01$ for random $u, v \in
\RR^N$. Two simultaneous branches of $\Mix(\alpha, T_1, \beta, T_2)$ coexist
without destructive interference; the ``choice'' collapses only at $\Project$
via the dominant activation axis, making $\Mix$ a sound replacement for
binary branching~(P4).
\end{remark}

\begin{theorem}[$\Project$ is a Functor]
\label{thm:functor}
Let $\mathcal{G} = (\mathrm{Pos}, \mathrm{Adm})$ be the \emph{geometric
category}: objects are positions $P \in \RR^N$; morphisms are admissible
transformations. Let $\mathcal{E}_X = (\mathrm{Type}_X, \mathrm{Map}_X)$ be the
\emph{executor category} of $\DimX$. Define:

\begin{description}[leftmargin=1.5em]
\item[\textbf{H1}] \emph{Embedding Coverage:} $\EmbX\colon
  \Ob(\mathcal{E}_X) \to \RR^N$ satisfies $\ProjX(\EmbX(t)) = t$ for all $t$,
  where $\ProjX(v) = \arg\min_{t} \norm{v - \EmbX(t)}_2$.
\item[\textbf{H2}] \emph{Subspace Preservation:} $T \in \mathrm{Adm}$ is
  $\DimX$-preserving if $\ProxX T \ProxX = \ProxX T$ (i.e., $T$ maps
  $\DimX$-active positions to $\DimX$-active positions).
\end{description}

Let $T_X \subseteq \mathrm{Adm}$ denote the $\DimX$-preserving transformations.
Define:
\[
  \Project_X\colon \mathcal{G}\big|_{T_X} \to \mathcal{E}_X
\]
\[
  \Project_X(P) = \ProjX(\ProxX V_P), \quad
  \Project_X(T) = \ProjX \circ (\ProxX T|_{\DimX}) \circ \EmbX.
\]
Under \textbf{H1} and \textbf{H2}, $\Project_X$ is a functor satisfying:
\begin{enumerate}[label=\textup{(\alph*)}]
\item \emph{Identity:} $\Project_X(\mathrm{id}_\mathcal{G}) = \mathrm{id}_{\mathcal{E}_X}$
\item \emph{Composition:} $\Project_X(T_2 \circ T_1) = \Project_X(T_2) \circ
  \Project_X(T_1)$ for all $T_1, T_2 \in T_X$
\end{enumerate}
\end{theorem}

\begin{proof}
\textbf{(a)} For any $t \in \Ob(\mathcal{E}_X)$: $\EmbX(t) \in \DimX$ by
design, so $\ProxX(\EmbX(t)) = \EmbX(t)$. Then $\Project_X(\mathrm{id})(t) =
\ProjX(\EmbX(t)) = t$ by~H1.

\textbf{(b)} Let $T_1, T_2 \in T_X$ and $t \in \Ob(\mathcal{E}_X)$; set
$v = \EmbX(t) \in \DimX$. By~H2 for $T_1$: $T_1(v) = w_1 \in \DimX$. By~H2
for $T_2$: $T_2(w_1) = w_2 \in \DimX$. Since $w_2 \in \DimX$, $\ProxX(w_2) =
w_2$. Therefore:
\[
  \Project_X(T_2 \circ T_1)(t) = \ProjX(\ProxX(T_2(T_1(v)))) = \ProjX(T_2(T_1(v))).
\]
For the right-hand side, let $t_1 = \Project_X(T_1)(t) = \ProjX(T_1(v))$.
Since $T_1(v) \in \DimX$, H1 gives $\EmbX(t_1) = T_1(v)$. Therefore:
\[
  \Project_X(T_2)(t_1) = \ProjX(T_2(\EmbX(t_1))) = \ProjX(T_2(T_1(v))).
\]
Both sides are equal for all $t$. \qed
\end{proof}

\begin{remark}
H2 holds by construction in Bootstrap~V0.3 and for transformations trained on
$\DimX$-specific data. In the approximate case---when $T$ leaks $\varepsilon$
into $\DimX^\perp$---Theorem~\ref{thm:functor} holds up to $O(\varepsilon)$.
\end{remark}

\begin{corollary}[Cross-Executor Portability]
\label{cor:portability}
For $\Pi = T_n \circ \cdots \circ T_1$ with all $T_i \in T_X$:
\[
  \Project_X(\Pi) = \Project_X(T_n) \circ \cdots \circ \Project_X(T_1).
\]
\end{corollary}

\begin{corollary}[Intersection as Pullback]
\label{cor:pullback}
$\Project(P, \mathsf{DIM}_A \cap \mathsf{DIM}_B)$ is the pullback of
$\Project(P, \mathsf{DIM}_A)$ and $\Project(P, \mathsf{DIM}_B)$ in
$\mathcal{E}_A \times \mathcal{E}_B$, provided $\mathsf{DIM}_A \cap
\mathsf{DIM}_B \neq \emptyset$ in $\RR^N$. This produces a live-bound
type---e.g., a Flutter widget synchronized with its SQL source---that has no
analogue in single-paradigm languages.
\end{corollary}

\begin{theorem}[Uniqueness of $\Fixpoint$]
\label{thm:fixpoint}
If $T$ is a contraction ($\norm{T(u)-T(v)} \leq k\norm{u-v}$, $k < 1$), then
$\exists!\ s^* \in \RR^N$ with $T(s^*) = s^*$, and $\Fixpoint(T, \varepsilon)$
converges to $s^*$ from any $s_0$.
\end{theorem}

\begin{proof}
Immediate from the Banach Fixed Point Theorem on $(\RR^N, \norm{\cdot}_2)$. \qed
\end{proof}

\paragraph{Denotational Semantics.}
$\llbracket\cdot\rrbracket \colon \mathit{Syntax} \to (\mathit{Env} \to
\mathit{Val})$ where $\mathit{Val} = \{T\colon \RR^N \to \RR^N \mid
\text{admissible}\}$:
\begin{align*}
\llbracket\Compose(T_1, T_2)\rrbracket\rho &= \llbracket T_2\rrbracket\rho \circ \llbracket T_1\rrbracket\rho \\
\llbracket\Mix(\alpha,T_1,\beta,T_2)\rrbracket\rho &= \alpha\llbracket T_1\rrbracket\rho + \beta\llbracket T_2\rrbracket\rho \\
\llbracket\Fixpoint(T,\varepsilon)\rrbracket\rho &= \lfp_\varepsilon(\llbracket T\rrbracket\rho) \\
\llbracket\Project(T,E)\rrbracket\rho &= F_E \circ \llbracket T\rrbracket\rho
\end{align*}
Functoriality of $\llbracket\mathsf{COMPILE}\rrbracket$ follows from
Theorem~\ref{thm:functor}.

\subsection{Illustrative Programs}

\begin{example}[Semantic Search with Live Interface]
\begin{lstlisting}[language=polydim,caption={Semantic search converging to a stable result projected simultaneously to Flutter and SQL.}]
position query : R^10000

transform attend_corpus :=
    ATTEND(W_Q, W_K_corpus, W_V_corpus)
transform refine :=
    MIX(0.8, attend_corpus,
        0.2, ATTEND(W_Q2, W_K2, W_V2))
transform search :=
    FIXPOINT(refine, 1e-4)

output ui   := RENDER(search, DIM_FLUTTER)
output db   := EXPORT(search, DIM_SQL)
-- Intersection type (Corollary 2):
output live := RENDER(search,
                 DIM_FLUTTER ∩ DIM_SQL)
\end{lstlisting}
\texttt{live} has a type that does not exist in any single-paradigm language: a
Flutter widget synchronized live with its SQL source, produced by the categorical
pullback of Corollary~\ref{cor:pullback}.
\end{example}

\begin{example}[Zero-Serialization AI-to-AI Delegation]
\begin{lstlisting}[language=polydim,caption={IA\_A delegates a task to IA\_B via the ALIGN protocol---no text, no JSON.}]
transform task :=
    COMPOSE(
        ATTEND(W_Q_task, W_K_ctx, W_V_ctx),
        MIX(0.9, COMPILE(
              ATTEND(W_Q2, W_K2, W_V2),
              DIM_RUST),
            0.1, RECUR(A, B, C)))
-- ALIGN transparent: runtime computes
-- T_B = M * task * M^+ before delivery
output delegation := COMPILE(task, DIM_RUST)
\end{lstlisting}
\emph{Comparison:} Text protocol (MCP/A2A) performs two lossy projections
($\RR^{12288} \to 100\,\text{tokens} \to \RR^{8192}$). POLYDIM~ALIGN performs
$T_B = M T_A M^\dagger$ in $O(d^2)$ with zero semantic loss.
\end{example}

%% ============================================================
\section{AI-to-AI Communication}
\label{sec:aitai}
%% ============================================================

\subsection{Cross-Attention as Native Communication}

Reading a POLYDIM message from $\mathrm{IA}_A$ is mathematically identical to
computing a cross-attention layer:
\[
  \mathrm{Communication}(s_B, s_A)
  = \mathrm{softmax}\!\left(\frac{Q_B K_A^\top}{\sqrt{d}}\right) V_A
\]
where $Q_B = s_B W_Q$, $K_A = s_A W_K$, $V_A = s_A W_V$. This is not a
metaphor but the exact mathematical operation occurring in multi-encoder
transformers.

\subsection{The ALIGN Protocol}

For $\mathrm{IA}_A \in \RR^{d_A}$ and $\mathrm{IA}_B \in \RR^{d_B}$:
\begin{enumerate}[label=\arabic*.,leftmargin=*]
\item \textbf{Pre-alignment:} Share codebook $\mathcal{C} = \{g_1,\ldots,g_k\}$
  of universal $\GeoID$s. Project: $A \in \RR^{d_A \times k}$, $B \in
  \RR^{d_B \times k}$.
\item \textbf{Procrustes:} $M^* = \arg\min_{M^\top M = I} \norm{MA - B}_F$.
  Solution: $\mathrm{SVD}(BA^\top) = U\Sigma V^\top$, $M = UV^\top$.
\item \textbf{Transformation:} $T_B = M T_A M^\dagger$ ($M^\dagger$ =
  Moore-Penrose pseudoinverse).
\item \textbf{Application:} $s_B' = T_B(s_B)$.
\end{enumerate}

\textbf{LoRA efficiency:} For $T_A = U_A V_A^\top$, we have $T_B =
(MU_A)(V_A M^\dagger)^\top$---the LoRA rank $r$ is preserved.

\textbf{HoTT-based $\GeoID$:} $[g] = \{s \in \RR^N \mid \exists \text{ admissible path } \gamma\colon g \rightsquigarrow s\}$. Two models identify the same object when their representations belong to the same homotopy class; the physical vector is a local chart, the class is the universal entity.

%% ============================================================
\section{Engineering}
\label{sec:engineering}
%% ============================================================

\paragraph{Binary format (.polydim V5, LoRA).}
The original V2 specification stored float32[$N{\times}N$] matrices: for $N =
10{,}000$, $400\,\mathrm{MB}$ per transform. V5 mandates $T = W_0 + UV^\top$
with $U, V \in \RR^{N \times r}$, yielding $2Nr \times 2\,\mathrm{B} \approx
2.56\,\mathrm{MB}$ for $r = 64$---a $\mathbf{156\times}$ reduction.

\paragraph{Sparse execution.}
The activation vector $A\colon D \to [0,1]$ acts as a semantic page table: if
$A(\mathsf{DIM}_i) < \tau$ (where $\tau = 0.5 + 2/(2\sqrt{N})$), the $i$-th
tensor block is skipped. Average complexity: $O(k_{\mathrm{active}} \cdot N
\cdot r)$ with $k_{\mathrm{active}} \ll 9$.

%% ============================================================
\section{Human Interface: Flutter Isomorphism}
\label{sec:flutter}
%% ============================================================

Flutter's widget tree admits a direct algebraic isomorphism to \textsc{Polydim}'s
composition algebra. The functor $\varphi\colon \mathcal{T} \to \mathcal{W}$:
\begin{align*}
\varphi(\Compose(T_1,T_2))    &\cong \texttt{Column([}\varphi(T_1), \varphi(T_2)\texttt{])} \\
\varphi(\Mix(\alpha,T_1,\beta,T_2)) &\cong \texttt{Stack([Opacity(}\alpha,\varphi(T_1)\texttt{),...}]} \\
\varphi(\Fixpoint(T,\varepsilon)) &\cong \texttt{StatefulWidget(tol:}\varepsilon\texttt{)}
\end{align*}

JS, React, and Java cannot serve as native executors: React's VDOM diffing has
no isomorphism to the composition algebra; Java's OOP model contradicts
geometric composition. Flutter is the unique web-deployed framework satisfying
the algebraic isomorphism.

When the user interacts with $\Project(\mathit{state}, \mathsf{DIM\_FLUTTER})$,
the action generates $\delta_v \in \RR^N$ injected directly into the geometric
state. Zero text serialization: 2 steps vs.\ 8 in current AI-UI pipelines.

%% ============================================================
\section{Related Work}
\label{sec:related}
%% ============================================================

\paragraph{Tensor frameworks.}
PyTorch~\citep{paszke2019pytorch}, JAX~\citep{bradbury2018jax}, and
TensorFlow~\citep{abadi2016tensorflow} use tensors as \emph{data} flowing
through sequential host-language programs. \textsc{Polydim} uses transformations
\emph{as} the program. The distinction is categorical: programming \emph{with}
tensors vs.\ programming \emph{as} tensors.

\paragraph{Agent communication.}
MCP~\citep{anthropic2024mcp} and A2A~\citep{google2025a2a} rely on JSON-RPC,
perpetuating the impedance mismatch. FIPA~\citep{fipa2002} grounds communication
in speech act theory requiring symbolic vocabulary agreement. POLYDIM eliminates
all three layers via geometric transformation transmission.

\paragraph{Categorical programming languages.}
Haskell~\citep{jones2003haskell}, Idris~\citep{brady2013idris}, and
Agda~\citep{norell2009agda} apply categorical semantics to type theory of
sequential programs. The closest antecedents are Categorical Deep
Learning~\citep{gavranovic2024categorical} and Backprop as
Functor~\citep{fong2019backprop}. \textsc{Polydim} extends both to the full
language-design level, with $\Project$ as the canonical functor.

\paragraph{Vector Symbolic Architectures.}
VSA research~\citep{plate1995,kanerva2009,frady2021} establishes algebraic
operations preserving semantic structure. \textsc{Polydim} extends VSA in two
directions: quasi-orthogonality founds $\Mix$ as sound if/else replacement, and
VSA binding underlies position binding and cross-model identity.

\paragraph{Compiler infrastructure.}
MLIR~\citep{lattner2021mlir} provides multi-level IR for heterogeneous ML
compilation. \textsc{Polydim} can be positioned as a \emph{Categorical IR}:
programs are morphisms in $\mathcal{G}$; compilation ($\Project$) is a functor.

\paragraph{Representation engineering.}
\citet{turner2023activation} and \citet{zou2023representation} demonstrate
targeted activation-space interventions. \textsc{Polydim} generalizes both:
$A\colon D \to [0,1]$ exposes activation weights as first-class program
parameters rather than implicit network behaviors.

%% ============================================================
\section{Discussion}
\label{sec:discussion}
%% ============================================================

\paragraph{What POLYDIM is not.}
\textsc{Polydim} is not a tensor computation framework (PyTorch programs compute
\emph{with} tensors; \textsc{Polydim} programs \emph{are} tensors), not a neural
architecture (transformers are the substrate), not a new attention mechanism
($\Attend$ is a pre-existing operation exposed as a primitive), and not a prompt
language (the text layer is bypassed entirely).

\paragraph{The assembly language parallel.}
Assembly did not invent the transistor; it formalized binary arithmetic
programming. \textsc{Polydim} does not invent the transformer; it formalizes
geometric computation programming. Unlike Assembly, \textsc{Polydim} is
model-agnostic via $\Align$---closer to LLVM IR than to x86 Assembly.

\paragraph{Open problems.}
\textbf{Q1} (Subspace emergence): How do the nine native subspaces emerge from
a real embedding model without manual definition? \textbf{Q2} (Distributed
consistency): Geometric-CRDTs where mutations converge via $\Mix$ and
Theorem~\ref{thm:fixpoint}. \textbf{Q3} (Turing completeness): Conjectured via
$\Fixpoint$ over contractive $T$ on compact $\RR^N$. \textbf{Q4} (Self-hosting):
Can a \textsc{Polydim} program generate \textsc{Polydim} programs?

%% ============================================================
\section{Conclusion}
\label{sec:conclusion}
%% ============================================================

We have presented \textsc{Polydim}, a programming language whose fundamental
unit is a geometric transformation $T\colon \RR^N \to \RR^N$, eliminating the
impedance mismatch between the multidimensional geometry of modern AI
architectures and the sequential text channels through which they currently
communicate. Four algebraic theorems establish the formal foundation.
Theorem~\ref{thm:functor}---$\Project$ is a functor---is the central
mathematical contribution: Corollary~\ref{cor:portability} guarantees
cross-executor portability, and Corollary~\ref{cor:pullback} establishes
intersection types as categorical pullbacks, producing live-bound types with no
analogue in single-paradigm languages. The Python bootstrap validates the
framework with $29/29$ tests and a $64.5\%$ semantic improvement. The Rust core
and LoRA binary format complete the engineering foundation. The assembly language
of the transformer era is here.

%% ============================================================
\bibliographystyle{plainnat}
\bibliography{polydim}
\end{document}

%%% ============================================================
%%% polydim.bib  (save as separate file)
%%% ============================================================
%
% @inproceedings{vaswani2017,
%   author    = {Vaswani, A. and Shazeer, N. and Parmar, N. and Uszkoreit, J.
%                and Jones, L. and Gomez, A.N. and Kaiser, L. and Polosukhin, I.},
%   title     = {Attention is All You Need},
%   booktitle = {Advances in Neural Information Processing Systems},
%   volume    = {30},
%   year      = {2017}
% }
% @article{gu2023mamba,
%   author  = {Gu, A. and Dao, T.},
%   title   = {Mamba: Linear-Time Sequence Modeling with Selective State Spaces},
%   journal = {arXiv:2312.00752},
%   year    = {2023}
% }
% @article{plate1995,
%   author  = {Plate, T.},
%   title   = {Holographic Reduced Representations},
%   journal = {IEEE Transactions on Neural Networks},
%   volume  = {6}, number = {3}, pages = {623--641},
%   year    = {1995}
% }
% @article{kanerva2009,
%   author  = {Kanerva, P.},
%   title   = {Hyperdimensional Computing},
%   journal = {Cognitive Computation},
%   volume  = {1}, number = {2}, pages = {139--159},
%   year    = {2009}
% }
% @article{elhage2022,
%   author  = {Elhage, N. and Hume, T. and Olsson, C. and others},
%   title   = {Toy Models of Superposition},
%   journal = {Transformer Circuits Thread},
%   year    = {2022}
% }
% @inproceedings{fong2019backprop,
%   author    = {Fong, B. and Spivak, D.I.},
%   title     = {Backprop as Functor},
%   booktitle = {LICS},
%   year      = {2019}
% }
% @article{gavranovic2024categorical,
%   author  = {Gavranovi{\'c}, B. and Lessard, P. and Dudzik, A. and others},
%   title   = {Categorical Deep Learning: An Algebraic Theory of Architectures},
%   journal = {arXiv:2402.15332},
%   year    = {2024}
% }
% @inproceedings{lattner2021mlir,
%   author    = {Lattner, C. and Amini, M. and Bondhugula, U. and others},
%   title     = {{MLIR}: Scaling Compiler Infrastructure},
%   booktitle = {CGO},
%   year      = {2021}
% }
% @inproceedings{paszke2019pytorch,
%   author    = {Paszke, A. and Gross, S. and Massa, F. and others},
%   title     = {{PyTorch}: An Imperative Style Deep Learning Library},
%   booktitle = {NeurIPS},
%   year      = {2019}
% }
% @software{bradbury2018jax,
%   author = {Bradbury, J. and Frostig, R. and Hawkins, P. and others},
%   title  = {{JAX}: Composable Transformations of {Python+NumPy}},
%   year   = {2018}
% }
% @inproceedings{abadi2016tensorflow,
%   author    = {Abadi, M. and Barham, P. and Chen, J. and others},
%   title     = {{TensorFlow}: A System for Large-Scale Machine Learning},
%   booktitle = {OSDI},
%   year      = {2016}
% }
% @article{frady2021,
%   author  = {Frady, E.P. and Kymn, D. and Sommer, F.T.},
%   title   = {Holographic Reduced Representations with {Fourier} Lift},
%   journal = {Cognitive Computation},
%   year    = {2021}
% }
% @article{moggi1991,
%   author  = {Moggi, E.},
%   title   = {Notions of Computation and Monads},
%   journal = {Information and Computation},
%   volume  = {93}, number = {1},
%   year    = {1991}
% }
% @book{jones2003haskell,
%   editor    = {Jones, S.P.},
%   title     = {{Haskell} 98 Language and Libraries},
%   publisher = {Cambridge University Press},
%   year      = {2003}
% }
% @article{brady2013idris,
%   author  = {Brady, E.},
%   title   = {{Idris}, a General-Purpose Dependently Typed Programming Language},
%   journal = {Journal of Functional Programming},
%   year    = {2013}
% }
% @inproceedings{norell2009agda,
%   author    = {Norell, U.},
%   title     = {Dependently Typed Programming in {Agda}},
%   booktitle = {TLDI},
%   year      = {2009}
% }
% @techreport{anthropic2024mcp,
%   author      = {{Anthropic}},
%   title       = {Model Context Protocol ({MCP})},
%   institution = {Anthropic},
%   year        = {2024}
% }
% @techreport{google2025a2a,
%   author      = {{Google}},
%   title       = {Agent-to-Agent Protocol ({A2A})},
%   institution = {Google},
%   year        = {2025}
% }
% @techreport{fipa2002,
%   author      = {{FIPA}},
%   title       = {{FIPA} Agent Communication Language},
%   institution = {Foundation for Intelligent Physical Agents},
%   year        = {2002}
% }
% @article{turner2023activation,
%   author  = {Turner, A. and Thiergart, L. and Udell, G. and others},
%   title   = {Activation Addition: Steering Language Models Without Optimization},
%   journal = {arXiv:2308.10248},
%   year    = {2023}
% }
% @article{zou2023representation,
%   author  = {Zou, A. and Phan, L. and Chen, S. and others},
%   title   = {Representation Engineering: A Top-Down Approach to {AI} Transparency},
%   journal = {arXiv:2310.01405},
%   year    = {2023}
% }
